\documentclass[11pt]{article}

\usepackage[margin=1in]{geometry}
\usepackage[T1]{fontenc}
\usepackage[utf8]{inputenc}
\usepackage{lmodern}
\usepackage{microtype}
\usepackage[round]{natbib}
\usepackage{hyperref}
\usepackage{enumitem}
\usepackage{booktabs}

\hypersetup{
  colorlinks=true,
  linkcolor=blue,
  citecolor=blue,
  urlcolor=blue
}

\setlist[itemize]{leftmargin=1.5em}
\setlist[enumerate]{leftmargin=1.5em}
\def\UrlBreaks{\do\/\do-\do\_\do\.\do\?\do\=\do\&}
\Urlmuskip=0mu plus 2mu\relax

\title{The Shovel-Ready Trap:\\ \large Project Pipelines and the Uneven Capacity to Do Regional Development\thanks{Draft date: 14 May 2026. Replication files, data manifests, claim ledgers, and reviewer-audit outputs are included in the public scriptor workflow.}}
\author{}
\date{May 2026}

\begin{document}
\maketitle

\begin{abstract}
Regional science has made major progress in explaining left-behind places, development traps, resilience, institutions, and the uneven outcomes of place-based policy. This paper adds a missing middle: the project pipeline through which places convert policy opportunity into bids, awards, procurements, contracts, and delivered interventions. I argue that competitive and projectified regional policy can create a shovel-ready trap. Places with stronger administrative capacity, fiscal slack, consultant access, planning bandwidth, and pre-existing project shelves are better able to turn funding windows into executable projects; places with high need may be filtered out before impact evaluation begins. I demonstrate the empirical object with a reproducible seed dataset for the UK Levelling Up Fund Round 2. The official assessment note records 529 bids seeking \pounds8.8 billion and 111 successful awards worth about \pounds2.09 billion; the committed parser recovers 111 award records from the official successful-bidders spreadsheet, implying a 21.0 percent bid-to-award conversion rate by count. The contribution is therefore conceptual, empirical, and infrastructural: it defines the project-conversion mechanism, shows that the first conversion stage is observable, and supplies a pipeline for extending the data to procurement, contracts, and delivery.
\end{abstract}

\input{sections/01_introduction}
\input{sections/02_literature_gap}
\input{sections/03_project_pipeline_concept}
\input{sections/04_data}
\input{sections/05_empirical_strategy}
\input{sections/06_expected_patterns}
\input{sections/07_discussion}
\input{sections/08_limits_audit}
\input{sections/09_conclusion}

\pagebreak
\bibliographystyle{plainnat}
\bibliography{references}

\end{document}
